% Minimal preamble for Open Logic Project sources (pinned commit 1e960beff9ed7835bf3e3f1335e21af3439cd107).
%
% OLP's own master preamble cannot run on the project host: TeX Live 2025 there lacks
% prooftree.sty, ntheorem.sty, extarrows.sty, mathdots.sty, isomath.sty and subfiles.sty.
% Following the project's own rule ("if the author preamble needs packages the toolchain lacks,
% do not invent macros -- use a minimal working preamble"), this file declares, verbatim in
% spirit and numbering behaviour, exactly the theorem-like environments OLP defines in
% sty/open-logic-book-envs.sty (thm/lem/prop/cor/defn/ex/prob/rem/note/case/axiom/conv),
% plus the shorthand accents OLP uses (`!!a`, `!!b`, `!!{...}`).
%
% Author text is untouched; only the packaging is ours.
\documentclass[10pt,a4paper]{article}

\usepackage{amsmath}
\usepackage{amssymb}
\usepackage{amsthm}
\usepackage{graphicx}
\usepackage{hyperref}
\usepackage{enumitem}
\usepackage[framemethod=default]{mdframed}

% ---- OLP shorthand accents (from sty/open-logic.sty) ----
% NB: \! is taken by LaTeX (negative thin space), so we only define the !!a / !!b families
% and the \olup helper used by them.
\newcommand*{\olup}[1]{\MakeUppercase{#1}}
\def\!!!a#1{\olup{#1}}
\def\!a#1{\olup{#1}}
\def\!b#1{\olup{#1}}

% ---- theorem-like environments, same family numbering as OLP ----
\theoremstyle{plain}
\newtheorem{thm}{Theorem}[section]
\newtheorem{lem}[thm]{Lemma}
\newtheorem{prop}[thm]{Proposition}
\newtheorem{cor}[thm]{Corollary}
\newtheorem{axiom}{Axiom}

\theoremstyle{definition}
\newtheorem{defn}[thm]{Definition}
\newtheorem{ex}[thm]{Example}
\newtheorem{prob}{Problem}[section]
\newtheorem{probd}{Problem}
\newtheorem{conv}[thm]{Convention}

\theoremstyle{remark}
\newtheorem{rem}[thm]{Remark}
\newtheorem{note}[thm]{Note}
\newtheorem{case}[thm]{Case}

% ---- environments OLP defines as wrappers ----
\newenvironment{probtag}[1]{\begin{prob}}{\end{prob}}
\newenvironment{explain}{}{}
\newenvironment{explain*}{}{}
\newenvironment{editorial}{}{}
\newenvironment{digress}{}{}
\newenvironment{digress*}{}{}
\newenvironment{history}{}{}
\newenvironment{reference}{}{}
\newcommand*{\olfileid}[3]{}
\newcommand*{\olsection}[1]{\section*{#1}}
\newcommand*{\olchapter}[3]{}
\newcommand*{\olpart}[2]{}
\newcommand*{\olimport}[1]{}
\newcommand*{\olinclude}[1]{}
\newcommand*{\olpath}[1]{}
\newcommand*{\ollangid}[1]{}
\newcommand*{\problemtag}[1]{}
\newcommand*{\problemauthor}[1]{}
\newcommand*{\ollabel}[1]{\label{#1}}
% \olref takes 0-3 OPTIONAL arguments plus one mandatory label:
%     \olref{lab}  |  \olref[from]{lab}  |  \olref[from][to]{lab}  |  \olref[a][b][c]{lab}
% A single-argument definition silently swallows "[sth]" and emits a broken reference
% (measured on the host: the text came out as "??iso]x"). The chain below consumes the
% brackets and cross-references only the mandatory argument.
\makeatletter
\newcommand{\olref}{\@ifnextchar[{\olref@i}{\olref@final}}
\def\olref@i[#1]{\@ifnextchar[{\olref@ii}{\olref@final}}
\def\olref@ii[#1]{\@ifnextchar[{\olref@iii}{\olref@final}}
\def\olref@iii[#1]{\olref@final}
\def\olref@final#1{\ref{#1}}
\makeatother
\newcommand*{\olnewpage}{\clearpage}

% ===================== OLP compatibility layer (1) =====================

\makeatletter

% ---------------- proof trees (OLP's prooftree is a bussproofs derivative) ----------------
\providecommand{\AxiomC}[1]{}
\providecommand{\UnaryInfC}[1]{}
\providecommand{\BinaryInfC}[1]{}
\providecommand{\TrinaryInfC}[1]{}
\providecommand{\QuaternaryInfC}[1]{}
\providecommand{\QuinaryInfC}[1]{}
\providecommand{\RightLabel}[1]{}
\providecommand{\LeftLabel}[1]{}
\providecommand{\doubleLine}{}
\providecommand{\noLine}{}
\providecommand{\DisplayProof}{}
\providecommand{\fCenter}{}
\newenvironment{prooftree}{}{}
\newenvironment{ltprooftree}{}{}
\newenvironment{oltableau}{}{}

% ---------------- logical notation ----------------
\providecommand{\lif}{\mathbin{\rightarrow}}
\providecommand{\lforall}{\mathord{\forall}}
\providecommand{\lexists}{\mathord{\exists}}
\providecommand{\ltrue}{\top}
\providecommand{\lfalse}{\bot}
\providecommand{\Ltrue}{\top}
\providecommand{\Lfalse}{\bot}
\providecommand{\Sequent}{\mathrel{\Rightarrow}}
\providecommand{\Proves}{\mathrel{\vdash}}
\providecommand{\Entails}{\mathrel{\models}}
\providecommand{\Subst}[3]{#1[#2/#3]}
\providecommand{\sFmla}[2]{#1(#2)}
\providecommand{\pFmla}[2]{#1(#2)}
\providecommand{\formula}{\relax}
\providecommand{\True}{\top}
\providecommand{\False}{\bot}
\providecommand{\Args}[1]{#1}
\providecommand{\Ts}{\ensuremath{\top}}
\providecommand{\Fi}{\ensuremath{\bot}}
\providecommand{\lfalsum}{\bot}

% ---------------- objects, structures, metatheory ----------------
\providecommand{\Obj}{}
\providecommand{\Struct}{\ensuremath{\mathfrak{A}}}
\providecommand{\Domain}{\ensuremath{|\mathfrak{A}|}}
\providecommand{\Assign}{\ensuremath{s}}
\providecommand{\Sat}[2]{\ensuremath{#1 \models #2}}
\providecommand{\mSat}[3]{\ensuremath{#1 \models #2[#3]}}
\providecommand{\mModel}[2]{\ensuremath{#1 \models #2}}
\providecommand{\Lang}{\ensuremath{\mathcal{L}}}
\providecommand{\Th}{\ensuremath{\mathrm{Th}}}
\providecommand{\Log}{\ensuremath{\mathrm{L}}}
\providecommand{\tuple}[1]{\langle #1 \rangle}
\providecommand{\num}[1]{\underline{#1}}
\providecommand{\cfind}[1]{\varphi_{#1}}
\providecommand{\Nat}{\ensuremath{\mathbb{N}}}
\providecommand{\Mult}{\ensuremath{\mathrm{Mult}}}
\providecommand{\Add}{\ensuremath{\mathrm{Add}}}
\providecommand{\Exp}{\ensuremath{\mathrm{Exp}}}
\providecommand{\Zero}{\ensuremath{0}}
\providecommand{\Succ}{\ensuremath{\mathrm{S}}}
\providecommand{\lambd}{\ensuremath{\lambda}}
\providecommand{\bredpar}{\mathrel{\rightarrow_\beta}}
\providecommand{\bcd}[1]{\ensuremath{\ulcorner #1 \urcorner}}
\providecommand{\Set}[1]{\ensuremath{\{#1\}}}
\providecommand{\indcase}[1]{\par\noindent\emph{Case #1}.\quad}
\providecommand{\indfrm}[1]{\par\noindent\emph{Form #1}.\quad}
\providecommand{\iftag}[1]{#1}
\providecommand{\tagitem}[1]{\item[#1]}
\providecommand{\Elim}[1]{}
\providecommand{\Intro}[1]{}
\providecommand{\LeftR}[1]{}
\providecommand{\RightR}[1]{}
\providecommand{\TAss}[1]{}
\providecommand{\TRule}[1]{}

\makeatother
\providecommand{\olpath}{.}

% ===================== OLP compatibility layer (2) =====================
% --- second compatibility block: notation and shorthands used by the frozen OLP units ---
\providecommand{\CutCS}{\mathrm{Cut}}
\providecommand{\LeftSubproofLabel}{\relax}
\providecommand{\RightSubproofLabel}{\relax}
\providecommand{\maeh}{\relax}
\providecommand{\cheight}{\relax}
\providecommand{\ande}{\mathrm{and}_E}
\providecommand{\andi}{\mathrm{and}_I}
\providecommand{\ore}{\mathrm{or}_E}
\providecommand{\ori}{\mathrm{or}_I}
\providecommand{\typeof}{\mathord{:}}
\providecommand{\pValue}[1]{#1}
\providecommand{\Sent}{\mathrm{Sent}}
\providecommand{\cardplus}{\mathbin{+}}
\providecommand{\cardtimes}{\mathbin{\cdot}}
\providecommand{\cardexpo}{\mathbin{\hat{\ }}}
\providecommand{\cardnless}{\not<}
\providecommand{\cardle}{\leq}
\providecommand{\funfromto}[3]{#1\colon #2 \to #3}
\providecommand{\funimage}[1]{\operatorname{im}(#1)}
\providecommand{\fdefined}{\mathord{\downarrow}}
\providecommand{\canonord}{\operatorname{ord}}
\providecommand{\ordtimes}{\mathbin{\cdot}}
\providecommand{\ordeq}[2]{#1 = #2}
\providecommand{\ordneq}[2]{#1 \neq #2}
\providecommand{\ordtype}[1]{\operatorname{ord}(#1)}
\providecommand{\ordsucc}{\mathord{+1}}
\providecommand{\ordunion}{\cup}
\providecommand{\onesphere}{S^1}
\providecommand{\rotationsgroup}{\mathrm{SO}(2)}
\providecommand{\equivrep}[1]{[#1]}
\providecommand{\rlexless}{\prec}
\providecommand{\ZFC}{\mathrm{ZFC}}
\providecommand{\ZF}{\mathrm{ZF}}
\providecommand{\Rat}{\mathbb{Q}}
\providecommand{\Real}{\mathbb{R}}
\providecommand{\PosInt}{\mathbb{Z}^{+}}
\providecommand{\Id}{\mathrm{id}}
\providecommand{\Bin}{\mathrm{Bin}}
\providecommand{\pto}{\rightharpoonup}
\providecommand{\dom}{\operatorname{dom}}
\providecommand{\ran}{\operatorname{ran}}
\providecommand{\len}{\operatorname{len}}
\providecommand{\concat}{\mathbin{{}^\frown}}
\providecommand{\card}[1]{\lvert #1 \rvert}
\providecommand{\Pow}{\mathcal{P}}
\providecommand{\Setabs}[2]{\{\, #1 \mid #2 \,\}}
\providecommand{\comp}{\circ}
\providecommand{\Char}{\chi}
\providecommand{\Prf}{\operatorname{Prf}}
\providecommand{\Prov}{\operatorname{Prov}}
\providecommand{\Proj}{\operatorname{proj}}
\providecommand{\Gn}{\operatorname{gn}}
\providecommand{\gn}{\operatorname{gn}}
\providecommand{\Value}[1]{\operatorname{val}(#1)}
\providecommand{\Deduce}{\mathrm{D}}
\providecommand{\Dual}{\mathrm{d}}
\providecommand{\Ax}{\mathrm{Ax}}
\providecommand{\pSat}[3]{#1 \models #2\,[#3]}
\providecommand{\pAssign}{\sigma}
\providecommand{\bforall}{\forall}
\providecommand{\bexists}{\exists}
\providecommand{\liff}{\mathbin{\leftrightarrow}}
\providecommand{\FV}{\operatorname{FV}}
\providecommand{\UnaryInf}{\ensuremath{\mid}}
\providecommand{\BinaryInf}{\ensuremath{\mid}}
\providecommand{\depth}{\operatorname{depth}}
\providecommand{\oliflabeldef}[3]{#1}
\providecommand{\TMblank}{\sqcup}
\providecommand{\TMstroke}{\mid}
\providecommand{\TMstay}{\mathrm{S}}
\providecommand{\TMright}{\mathrm{R}}
\providecommand{\TMleft}{\mathrm{L}}
\providecommand{\TMendtape}{\triangleright}
\providecommand{\umin}{\mu}
\providecommand{\Sum}{\sum}
\providecommand{\citep}[1]{\cite{#1}}
\providecommand{\citet}[1]{\cite{#1}}
\providecommand{\citeauthor}[1]{\cite{#1}}
\providecommand{\citeyear}[1]{\cite{#1}}

% ===================== OLP compatibility layer (3) =====================
% --- third compatibility block: remaining notation seen in the frozen OLP units ---
\providecommand{\iso}{\cong}
\providecommand{\nszero}{\mathrm{0}}
\providecommand{\nssucc}{\mathrm{S}}
\providecommand{\nsplus}{\mathbin{+}}
\providecommand{\nssless}{<}
\providecommand{\nsless}{<}
\providecommand{\PL}{\mathrm{PL}}
\providecommand{\Nec}{\mathord{\Box}}
\providecommand{\RK}{\mathrm{RK}}
\providecommand{\ST}{\mathrm{ST}}
\providecommand{\SSubst}[3]{#1[#2/#3]}
\providecommand{\subst}[3]{#1[#2/#3]}
\providecommand{\tagfalse}{\mathrm{F}}
\providecommand{\cutr}{\mathrm{Cut}}
\providecommand{\cardneq}{\neq}
\providecommand{\issucc}[1]{\operatorname{succ}(#1)}
\providecommand{\tup}[1]{\langle #1 \rangle}
\providecommand{\Var}{\mathrm{Var}}
\providecommand{\Bew}{\operatorname{Bew}}
\providecommand{\Con}{\operatorname{Con}}
\providecommand{\Cons}{\operatorname{Cons}}
\providecommand{\pDeduce}{\mathrm{D}}
\providecommand{\psat}{\models}
\providecommand{\den}[1]{\llbracket #1 \rrbracket}
\providecommand{\Sem}[2]{\llbracket #1 \rrbracket_{#2}}
\providecommand{\DomainOf}[1]{\lvert #1 \rvert}
\providecommand{\Form}{\mathrm{Form}}
\providecommand{\Terms}{\mathrm{Term}}
\providecommand{\Atom}{\mathrm{Atom}}
\providecommand{\FVars}[1]{\operatorname{FV}(#1)}
\providecommand{\fv}[1]{\operatorname{fv}(#1)}
\providecommand{\tuplep}[1]{\langle #1 \rangle}
\providecommand{\restrict}[2]{#1 \upharpoonright #2}
\providecommand{\sto}{\rightarrow}
\providecommand{\FmlaSet}[1]{\mathrm{Fmla}_{#1}}
\providecommand{\Thm}{\mathrm{Thm}}
\providecommand{\Eval}{\operatorname{val}}
\providecommand{\blank}{\sqcup}
\providecommand{\blankTape}{\sqcup}
\providecommand{\Consis}{\operatorname{Con}}
\providecommand{\drs}{\mathrm{d}}
\providecommand{\VRule}{\mid}
\providecommand{\Bet}{\mathrm{B}}
\providecommand{\Seg}[1]{\mathrm{Seg}_{#1}}
\providecommand{\Enums}{\mathrm{En}}
\providecommand{\Z}{\mathbb{Z}}
\providecommand{\N}{\mathbb{N}}
\providecommand{\abs}[1]{\lvert #1 \rvert}
\providecommand{\floor}[1]{\lfloor #1 \rfloor}
\providecommand{\ceil}[1]{\lceil #1 \rceil}

% ===================== OLP compatibility layer (4) =====================
% --- fourth compatibility block: harvested from the per-candidate compile logs ---
\providecommand{\beredpar}{\mathrel{\rightarrow_\beta}}
\makeatletter
\def\bered{\@ifnextchar[{\bered@i}{\mathrel{\rightarrow_\beta}}}
\def\bered@i[#1]{\mathrel{\rightarrow_\beta}}
\def\bred{\@ifnextchar[{\bred@i}{\mathrel{\rightarrow_\beta}}}
\def\bred@i[#1]{\mathrel{\rightarrow_\beta}}
\makeatother
\providecommand{\beredone}{\mathrel{=_\beta}}
\providecommand{\bredone}{\mathrel{=_\beta}}
\providecommand{\xredone}{\mathrel{=_\beta}}
\providecommand{\aconvone}{\mathrel{=_\alpha}}
\providecommand{\cardless}{<}
\providecommand{\tagrefs}[1]{#1}
\providecommand{\ordplus}{\mathbin{+}}
\providecommand{\funrestrictionto}[2]{#1 \upharpoonright #2}
\providecommand{\fundefined}{\mathord{\uparrow}}
\providecommand{\Trm}{\mathrm{Trm}}
\providecommand{\OProv}{\operatorname{Prov}}
\providecommand{\OPrf}{\operatorname{Prf}}
\providecommand{\aeq}{\mathrel{=}}
\providecommand{\mClass}[1]{[#1]}
\providecommand{\equivclass}[1]{[#1]}
\providecommand{\pheight}{\operatorname{ht}}
\providecommand{\phy}{\operatorname{ht}}
\providecommand{\Cut}{\mathrm{Cut}}
\providecommand{\QuantRank}{\operatorname{qr}}
\providecommand{\elemequiv}{\mathrel{\equiv}}
\providecommand{\PIso}{\mathrm{PIso}}
\providecommand{\Theory}{\mathrm{Th}}
\providecommand{\supstrict}{\subsetneq}
\providecommand{\trcl}{\operatorname{trcl}}
\providecommand{\cref}[1]{\ref{#1}}
\providecommand{\Discharge}{\mathrm{Dis}}
\providecommand{\Weakening}{\mathrm{Wk}}
\providecommand{\ZFminus}{\mathrm{ZF}^{-}}
\makeatletter
\renewcommand{\Dual}{\@ifnextchar[{\Dual@i}{\ifmmode\mathrm{d}\else$\mathrm{d}$\fi}}
\def\Dual@i[#1]{\ifmmode\mathrm{d}[#1]\else$\mathrm{d}[#1]$\fi}
\makeatother
\providecommand{\RightR}{\mathrm{R}}
\providecommand{\LeftR}{\mathrm{L}}
\providecommand{\Intro}{\mathrm{I}}
\providecommand{\Axiom}{\mathrm{Ax}}
\providecommand{\formula}[1]{#1}
\providecommand{\stared}{\ast}
\providecommand{\oset}[2]{\{\, #1 \mid #2 \,\}}
\providecommand{\setof}[2]{\{\, #1 \mid #2 \,\}}
\providecommand{\pow}[1]{\mathcal{P}(#1)}
\providecommand{\ordpair}[2]{\langle #1, #2 \rangle}
\providecommand{\ordtriple}[3]{\langle #1, #2, #3 \rangle}
\providecommand{\union}{\cup}
\providecommand{\inter}{\cap}
\providecommand{\seq}[1]{\langle #1 \rangle}
\providecommand{\SuchThat}{\mid}
\providecommand{\where}{\mid}
\providecommand{\TVar}{\mathrm{TVar}}
\providecommand{\Bool}{\mathrm{Bool}}
\providecommand{\Val}{\operatorname{Val}}
\providecommand{\TL}[1]{\mathrm{TL}}
\providecommand{\xstared}{\ast}

% ===================== OLP compatibility layer (5) =====================
% --- fifth compatibility block: rescues the quarantined OLP candidates ---
\providecommand{\ident}{\operatorname{id}}
\providecommand{\setrank}{\operatorname{rank}}
\providecommand{\cardfont}{\mathfrak{c}}
\providecommand{\Complement}{\overline}
\providecommand{\QR}{\mathrm{QR}}
\providecommand{\becd}{\ensuremath{\ulcorner}}
\providecommand{\disjointsum}{\mathbin{\oplus}}
\providecommand{\cardeq}[2]{#1 = #2}
\providecommand{\xred}{\mathrel{\rightarrow}}
% --- sixth compatibility block: names reported by real compile logs ---
\providecommand{\ORefut}{\operatorname{Refut}}
\providecommand{\ORProv}{\operatorname{Prov}}
\providecommand{\cutrank}{\operatorname{cr}}
\providecommand{\Refut}{\operatorname{Refut}}
\providecommand{\DischargeRule}{\mathrm{Dis}}
\providecommand{\Contraction}{\mathrm{C}}
\providecommand{\bottomAlignProof}{\relax}
\providecommand{\FalseCl}{\bot}
\providecommand{\red}{\mathrel{\rightarrow}}
\providecommand{\varAssign}{\sigma}
\providecommand{\defiff}{\mathrel{\leftrightarrow}}
\providecommand{\FalseInt}{\bot}
\providecommand{\fn}{\operatorname{fn}}
\providecommand{\dotted}{\cdot}

% --- tag-macro semantic fix (2026-10-13, final): \iftag CHOOSES a branch and the OLP proof body is
% always the SECOND group (\iftag{tag}{body}, {tag}{body}{c}, {tag}{body}{c}{d}); a definition that
% prints the tag literal drops the proof body. Bodies contain blank lines, hence \long.
\makeatletter
\providecommand{\iftag}{}
\long\def\iftag#1#2{\@ifnextchar\bgroup{\iftag@three{#2}}{\iftag@stop{#2}}}
\long\def\iftag@three#1#2{\iftag@stop{#1}\@ifnextchar\bgroup{\iftag@eat}{\relax}}
\long\def\iftag@eat#1{}
\long\def\iftag@stop#1{#1}
\long\def\oliflabeldef#1#2{\@ifnextchar\bgroup{\@gobble}{\iftag@stop{#2}}}
\providecommand{\tagrefs}[1]{}
\providecommand{\settag}[1]{}
\makeatother
\begin{document}

% ==== OLP file: content/set-theory/spine/foundation.tex ====


\olfileid{sth}{spine}{foundation}

\olsection{Foundation}

We are only \emph{almost} done---and not \emph{quite} finished---because nothing in
$\ZFminus$ guarantees that \emph{every} set is in some $V_\alpha$,
i.e., that every set is formed at some stage. 

Now, there is a fairly straightforward (mathematical) sense in which
we don't \emph{care} whether there are sets outside the hierarchy. (If
there are any there, we can simply ignore them.) But we have motivated
our \emph{concept} of set with the thought that every set is formed at
some stage (see \stageshier{} in \olref[z][story]{sec}). So
we will want to preclude the possibility of sets which fall outside of
the hierarchy. Accordingly, we must add a new axiom, which ensures
that every set occurs somewhere in the hierarchy. 

Since the $V_\alpha$s are our stages, we might simply consider adding
the following as an axiom:

\begin{defish}
\emph{Regularity.} $\forall A \exists \alpha\, A \subseteq V_\alpha$
\end{defish}

This would be a perfectly reasonable approach. However, for reasons
that will be explained in the next section, we will instead adopt an
alternative axiom:

\begin{axiom}[Foundation]
$(\forall A \neq \emptyset)(\exists B \in A)A \cap B = \emptyset$.
\end{axiom}

With some effort, we can show (in $\ZFminus$) that Foundation entails Regularity:
\begin{defn}
For each set $A$, let:
\begin{align*}
	\text{cl}_0(A) &= A,\\
	\text{cl}_{n+1}(A) &= \bigcup \text{cl}_n(A),\\
	\text{trcl}(A) &= \bigcup_{n < \omega} \text{cl}_{n}(A).
\end{align*}
We call $\text{trcl}(A)$ the \emph{transitive closure} of $A$.
\end{defn}
\noindent
The name ``transitive closure'' is apt:
\begin{prop}\ollabel{subsetoftrcl}
$A \subseteq \trcl{A}$ and $\trcl{A}$ is a transitive set. 
\end{prop}

\begin{proof}
Evidently $A = \text{cl}_0(A) \subseteq \text{trcl}(A)$. And if $x
\in b \in \trcl{A}$,  then $b \in \text{cl}_n(A)$ for some $n$, so $x
\in \text{cl}_{n+1}(A) \subseteq \text{trcl}(A)$. 
\end{proof}

\begin{lem}\ollabel{lem:TransitiveWellFounded}
If $A$ is a transitive set, then there is some $\alpha$ such that $A
\subseteq V_\alpha$.
\end{lem}

\begin{proof}
Recalling the definition of ``$\supstrict(X)$'' from
\olref[ordinals][opps]{defsupstrict}, define two sets:
\begin{align*}
	D  &= \Setabs{x \in A}{\forall \delta\ x \nsubseteq V_\delta}\\
	\alpha &= \supstrict\Setabs{\delta}{(\exists x \in A)
	(x \subseteq V_\delta \land (\forall \gamma \in \delta)x \nsubseteq V_\gamma)}
\end{align*}
Suppose $D = \emptyset$. So if $x \in A$, then there is some $\delta$ such that $x \subseteq V_\delta$ and, by the well-ordering of the ordinals, $(\forall \gamma \in \delta)x \nsubseteq V_\gamma$; hence $\delta \in \alpha$ and so $x \in V_\alpha$ by
\olref[spine][Valphabasic]{Valphabasicprops}. Hence $A \subseteq
V_{\alpha}$, as required. 

So it suffices to show that $D = \emptyset$. For reductio, suppose
otherwise. By Foundation, there is some $B \in D \subseteq A$ such that $D \cap B
= \emptyset$. If $x \in B$ then $x \in A$, since $A$ is transitive,
and since $x \notin D$, it follows that $\exists \delta\ x \subseteq
V_\delta$. So now let
\[
	\beta = \supstrict\Setabs{\delta}{(\exists x \in b)(x \subseteq V_\delta \land (\forall \gamma < \delta)x \nsubseteq V_\gamma)}.
\]
As before, $B \subseteq V_\beta$, contradicting the claim that $B \in
D$.	
\end{proof}

\begin{thm}\ollabel{zfentailsregularity}
Regularity holds.
\end{thm}

\begin{proof}
Fix $A$; now $A \subseteq \trcl{A}$ by \olref{subsetoftrcl}, which is transitive. So there is some $\alpha$ such that $A \subseteq \trcl{A} \subseteq V_\alpha$ by \olref{lem:TransitiveWellFounded}
\end{proof}

These results show that $\ZFminus$ proves the conditional
$\emph{Foundation}\Rightarrow\emph{Regularity}$. In
\olref[spine][rank]{zfminusregularityfoundation}, we will show that
$\ZFminus$ proves $\emph{Regularity}\Rightarrow\emph{Foundation}$. As
such, Foundation and Regularity are \emph{equivalent} (modulo
$\ZFminus$). But this means that, given $\ZFminus$, we can justify
Foundation by noting that it is equivalent to Regularity. And we can
justify Regularity immediately on the basis of \stageshier{}. 

% ==== OLP file: content/set-theory/spine/idea.tex ====


\olfileid{sth}{spine}{valpha}

\olsection{Defining the Stages as the $V_\alpha$s}

In \olref[sth][ordinals][]{chap}, we defined well-orderings and the
(von Neumann) ordinals. In this chapter, we will use these to
characterise the hierarchy of sets \emph{itself}. To do this, recall
that in \olref[sth][ordinals][opps]{sec}, we defined the idea of
successor and limit ordinals. We use these ideas in following
definition:

\begin{defn}\ollabel{defValphas}
	\begin{align*}
	V_\emptyset &\defis \emptyset\\
	V_{\ordsucc{\alpha}} &\defis \Pow{V_\alpha} & & 
	\text{for any ordinal }\alpha\\
	V_{\alpha} &\defis \bigcup_{\gamma < \alpha} V_\gamma & & 
	\text{when }\alpha\text{ is a limit ordinal}
\end{align*}
\end{defn}
\noindent
This will be a definition by \emph{transfinite recursion} on the
ordinals. In this regard, we should compare this with recursive
definitions of functions on the natural numbers.\footnote{Cf.\ the
definitions of addition, multiplication, and exponentiation in
\olref[sfr][infinite][dedekind]{sec}.} As when dealing with natural
numbers, one defines a base case and successor cases; but when dealing
with ordinals, we also need to describe the behaviour of \emph{limit}
cases. 

This definition of the $V_\alpha$s will be an important milestone. We
have informally motivated our hierarchy of sets as forming sets by
\emph{stages}. The $V_\alpha$s are, in effect, just those stages.
Importantly, though, this is an \emph{internal} characterisation of
the stages. Rather than suggesting a possible \emph{model} of the
theory, we will have defined the stages \emph{within} our set theory.


% ==== OLP file: content/set-theory/spine/rank.tex ====


\olfileid{sth}{spine}{rank}
\olsection{Rank}

Now that we have defined the stages as the $V_\alpha$'s, and we know
that every set is a subset of some stage, we can define the
\emph{rank} of a set. Intuitively, the rank of $A$ is the first moment
at which $A$ is formed. More precisely:

\begin{defn}\ollabel{defnsetrank}
For each set $A$, $\setrank{A}$ is the least ordinal $\alpha$ such that $A
\subseteq V_\alpha$.
\end{defn}
\begin{prop}\ollabel{ranksexist}
	$\setrank{A}$ exists, for any $A$.
\end{prop}
\begin{proof}
	Left as an exercise.
\end{proof}
\begin{prob}
	Prove \olref[sth][spine][rank]{ranksexist}.	
\end{prob}
\noindent 
The well-ordering of ranks allows us to prove some important results:
\begin{prop}\ollabel{valphalowerrank}
For any ordinal $\alpha$, $V_\alpha = \Setabs{x}{\setrank{x} \in \alpha}$.
\end{prop}

\begin{proof}
If $\setrank{x} \in \alpha$ then $x \subseteq V_{\setrank{x}} \in
V_\alpha$, so $x \in V_\alpha$ as $V_\alpha$ is potent (invoking
\olref[Valphabasic]{Valphabasicprops} multiple times). Conversely, if $x \in V_\alpha$ then $x \subseteq V_\alpha$, so $\setrank{x} \leq \alpha$; now a simple transfinite induction shows that $x \notin V_\alpha$. 
\end{proof}
\begin{prob}
	Complete the simple transfinite induction mentioned in \olref[sth][spine][rank]{valphalowerrank}.
\end{prob}

\begin{prop}\ollabel{rankmemberslower}
If $B \in A$, then $\setrank{B} \in \setrank{A}$.
\end{prop}

\begin{proof}
$A \subseteq V_{\setrank{A}} = \Setabs{x}{\setrank{x} \in
\setrank{A}}$ by \olref{valphalowerrank}.
\end{proof}
\noindent
Using this fact, we can establish a result which allows us to prove
things about \emph{all sets} by a form of induction:

\begin{thm}[$\in$-Induction Scheme] 
For any formula $\phi$:
\[
	\forall A((\forall x \in A)\phi(x) \lif \phi(A)) \lif \forall A \phi(A).
\]
\end{thm}

\begin{proof}
We will prove the contrapositive. So, suppose $\lnot \forall A
\phi(A)$. By Transfinite Induction
(\olref[ordinals][basic]{ordinductionschema}), there
is some non-$\phi$ of least possible rank; i.e.\ some $A$
such that $\lnot \phi(A)$ and $\forall x(\setrank{x} \in \setrank{A}
\lif \phi(x))$. Now if $x \in A$ then $\setrank{x} \in
\setrank{A}$, by \olref{rankmemberslower}, so that $\phi(x)$; i.e.\ $(\forall x \in
A)\phi(x) \land \lnot \phi(A)$.
\end{proof}\noindent Here is an informal way to gloss this powerful
result. Say that $\phi$ is \emph{hereditary} iff whenever every
\olup{element} of a set is $\phi$, the set itself is $\phi$. Then
$\in$-Induction tells you the following: if $\phi$ is hereditary,
every set is $\phi$.

To wrap up the discussion of ranks (for now), we'll prove a few
claims which we have foreshadowed a few times. 

\begin{prop}\ollabel{ranksupstrict}
$\setrank{A} = \supstrict_{x \in A}\setrank{x}$.
\end{prop}

\begin{proof}
Let $\alpha = \supstrict_{x \in A}\setrank{x}$. By
\olref{rankmemberslower}, $\alpha \leq \setrank{A}$. But if $x \in A$
then $\setrank{x} \in \alpha$, so that $x \in V_\alpha$ by \olref{valphalowerrank}, and hence $A
\subseteq V_\alpha$, i.e., $\setrank{A} \leq \alpha$. Hence
$\setrank{A} = \alpha$.
\end{proof}

\begin{cor}\ollabel{ordsetrankalpha}
For any ordinal $\alpha$, $\setrank{\alpha} = \alpha$.
\end{cor}

\begin{proof}
Suppose for transfinite induction that $\setrank{\beta} = \beta$ for
all $\beta \in \alpha$. Now $\setrank{\alpha} = \supstrict_{\beta \in
\alpha}\setrank{\beta} = \supstrict_{\beta \in \alpha}\beta = \alpha$
by \olref{ranksupstrict}.
%	First note that $\setrank{\alpha} \neq \beta$ for any $\beta \in
%	\alpha$. For otherwise we would have $\beta = \setrank{\beta} \in
%	\setrank{\alpha} = \beta$ by \olref{rankmemberslower}, i.e.,
%	$\beta \in \beta$, a contradiction. 
%
%	Now note that $\alpha \subseteq V_\alpha$. For $\alpha =
%	\Setabs{\beta}{\beta \in \alpha}$ by
%	\olref[sfr][ordinals][basic]{ordissetofsmallerord}. And if
%	$\beta \in \alpha$ then $\beta \subseteq V_{\setrank{\beta}} =
%	V_\beta \in V_\alpha$, hence $\beta \in V_\alpha$, by
%	\olref[sfr][spine][Valphabasic]{Valphabasicprops} twice. 
\end{proof}

Finally, here is a quick proof of the result promised at the end of
\olref[foundation]{sec}, that $\ZFminus$ proves the conditional
$\emph{Regularity} \Rightarrow \emph{Foundation}$. (Note that the
notion of ``rank'' and \olref{rankmemberslower} are available for use
in this proof since---as mentioned at the start of this section---they
can be presented using $\ZFminus + \text{Regularity}$.)

\begin{prop}[working in $\ZFminus + \text{Regularity}$]\ollabel{zfminusregularityfoundation} Foundation holds.
\end{prop}

\begin{proof}
Fix $A \neq \emptyset$, and some $B \in A$ of least possible rank. If
$c \in B$ then $\setrank{c} \in \setrank{B}$ by
\olref{rankmemberslower}, so that $c \notin A$ by choice of $B$.
\end{proof}

% ==== OLP file: content/set-theory/spine/recursion.tex ====


\olfileid{sth}{spine}{recursion}
\olsection{The Transfinite Recursion Theorem(s)}

The first thing we must do, though, is confirm that
\olref[valpha]{defValphas} is a successful definition. More generally,
we need to prove that any attempt to offer a transfinite by
(transfinite) recursion will succeed. That is the aim of this section.

\emph{Warning: this is tricky material}. The overarching moral,
though, is quite simple: Transfinite Induction plus Replacement
guarantee the legitimacy of (several versions of) transfinite
recursion.\footnote{A reminder: all formulas and terms can have parameters (unless explicitly stated otherwise).}

\begin{defn}
	Let $\tau(x)$ be a term; let $f$ be a function; let $\alpha$ be an ordinal. We say that $f$ is an \emph{$\alpha$-approximation} for $\tau$ iff both $\dom{f} = \alpha$ and $(\forall \beta \in \alpha)f(\beta) = \tau(\funrestrictionto{f}{\beta})$.
\end{defn}

\begin{lem}[Bounded Recursion]\ollabel{transrecursionfun}
For any term $\tau(x)$ and any ordinal $\alpha$, there is a unique $\alpha$-approximation for $\tau$.
\end{lem}
\begin{proof}
We will show that, for any $\gamma \leq \alpha$, there is a unique
$\gamma$-approximation. 

We first establish uniqueness. Let $g$ and $h$ (respectively) be $\gamma$- and $\delta$-approximations. A transfinite induction on their arguments shows that
$g(\beta) = h(\beta)$ for any $\beta \in \dom{g} \cap \dom{h} =
\gamma \cap \delta = \min(\gamma, \delta)$. So our approximations are unique (if they exist), and agree on all values.

To establish existence, we now use a simple transfinite induction
(\olref[ordinals][opps]{simpletransrecursion}) on ordinals
$\delta \leq \alpha$. 

The empty function is trivially an $\emptyset$-approximation. 

If $g$ is a $\gamma$-approximation, then $g \cup
\{\tuple{\gamma, \tau(g)}\}$ is a $\ordsucc{\gamma}$-approximation.

If $\gamma$ is a limit ordinal and $g_\delta$ is a $\delta$-approximation for all $\delta < \gamma$, let $g = \bigcup_{\delta \in \gamma} g_\delta$. This
is a function, since our various $g_\delta$s agree on all values. And
if $\delta \in \gamma$ then $g(\delta) = g_{\ordsucc{\delta}}(\delta) =
\tau(\funrestrictionto{g_{\ordsucc{\delta}}}{\delta}) =
\tau(\funrestrictionto{g}{\delta})$.

This completes the proof by transfinite induction.
\end{proof}

If we allow ourselves to define a \emph{term} rather than a function,
then we can remove the bound $\alpha$ from the previous result. In
the statement and proof of the following result, when $\sigma$ is a term, we
let $\funrestrictionto{\sigma}{\alpha} = \Setabs{\tuple{\beta,
\sigma(\beta)}}{\beta \in \alpha}$.

\begin{thm}[General Recursion]\ollabel{transrecursionschema}
For any term $\tau(x)$, we can explicitly define a term
$\sigma(x)$, such that
$\sigma(\alpha) = \tau(\funrestrictionto{\sigma}{\alpha})$ for any
ordinal~$\alpha$.
\end{thm}

\begin{proof}
For each $\alpha$, by \olref{transrecursionfun} there is a unique $\alpha$-approximation, $f_\alpha$, for~$\tau$. Define $\sigma(\alpha)$ as $f_{\ordsucc{\alpha}}(\alpha)$. Now:
	\begin{align*}
		\sigma(\alpha) &= 
		f_{\ordsucc{\alpha}}(\alpha) \\&= 
		\tau(\funrestrictionto{f_{\ordsucc{\alpha}}}{\alpha}) \\&= 
		\tau(\Setabs{\tuple{\beta, f_{\ordsucc{\alpha}}(\beta)}}{\beta \in \alpha}) \\&= 
		\tau(\Setabs{\tuple{\beta, f_{\ordsucc{\beta}}(\beta)}}{\beta \in \alpha})\\&=
		\tau(\funrestrictionto{\sigma}{\alpha})
	\end{align*}
noting that $f_{\ordsucc{\beta}}(\beta) = f_{\ordsucc{\alpha}}(\beta)$ for all $\beta < \alpha$, as in \olref{transrecursionfun}. 
\end{proof}
\noindent 
Note that \olref{transrecursionschema} is a \emph{schema}. Crucially, we cannot
expect $\sigma$ to define a function, i.e., a certain kind of
\emph{set}, since then $\dom{\sigma}$ would be the set of all
ordinals, contradicting the Burali-Forti Paradox
(\olref[ordinals][basic]{buraliforti}).

It still remains to show, though, that \olref{transrecursionschema}
vindicates our definition of the $V_\alpha$s. This may not be
immediately obvious; but it will become apparent with a last, simple,
version of transfinite recursion.

\begin{thm}[Simple Recursion]\ollabel{simplerecursionschema} 
For any terms $\tau(x)$ and $\theta(x)$ and any set $A$, we can
explicitly define a term $\sigma(x)$ such that:
\begin{align*}
	\sigma(\emptyset) &= A\\
	\sigma(\ordsucc{\alpha}) &= \tau(\sigma(\alpha)) &&
		\text{for any ordinal }\alpha\\
	\sigma(\alpha) &= \theta(\ran{\funrestrictionto{\sigma}{\alpha}})&&
	\text{when }\alpha\text{ is a limit ordinal}
\end{align*}
\end{thm}

\begin{proof}
We start by defining a term, $\xi(x)$, as follows: 
%t $\xi(x) = A$ if $x$ is not a function whose domain is an ordinal;
%otherwise let:
\[
	\xi(x) = 
	\begin{cases}
			A &\text{if $x$ is not a function whose}\\
			&\hspace{1em}\text{domain is an ordinal; otherwise:}\\
			\tau(x(\alpha)) & \text{if $\dom{x} = \ordsucc{\alpha}$}\\
			\theta(\ran{x}) & \text{if $\dom{x}$ is a limit ordinal}
	\end{cases}
\]
By \olref{transrecursionschema}, there is a term $\sigma(x)$ such that
$\sigma(\alpha) = \xi(\funrestrictionto{\sigma}{\alpha})$ for every
ordinal $\alpha$; moreover, $\funrestrictionto{\sigma}{\alpha}$ is a
function with domain $\alpha$. We show that $\sigma$ has the required
properties, by simple transfinite induction
(\olref[ordinals][opps]{simpletransrecursion}). 

First, $\sigma(\emptyset) = \xi(\emptyset) = A$. 

Next, $\sigma(\ordsucc{\alpha}) = \xi(\funrestrictionto{\sigma}{\ordsucc{\alpha}}) = \tau(\funrestrictionto{\sigma}{\ordsucc{\alpha}}(\alpha)) = \tau(\sigma(\alpha))$.

Last, $\sigma(\alpha) = \xi(\funrestrictionto{\sigma}{\alpha}) = \theta(\ran{\funrestrictionto{\sigma}{\alpha}})$, when $\alpha$ is a limit.
\end{proof}
\noindent
Now, to vindicate \olref[valpha]{defValphas}, just take $A
= \emptyset$ and $\tau(x) = \Pow{x}$ and $\theta(x) = \bigcup x$. At long last, this vindicates the definition of the $V_\alpha$s!{}



% ==== OLP file: content/set-theory/spine/spine.tex ====
% Part: set-theory
% Chapter: spine



\olchapter{sth}{spine}{Stages and Ranks}

\olimport{idea}
\olimport{recursion}
\olimport{stagesbasics}
\olimport{foundation}
\olimport{zf}
\olimport{rank}

\OLEndChapterHook

% ==== OLP file: content/set-theory/spine/stagesbasics.tex ====


\olfileid{sth}{spine}{Valphabasic}
\olsection{Basic Properties of Stages}

To bring out the foundational importance of the definition of the
$V_\alpha$s, we will present a few basic results about them. We start with a definition:\footnote{There's no standard terminology for ``potent''; this is the name used by \citet{ButtonLT1}.}
\begin{defn}
	The set $A$ is \emph{potent} iff $\forall x((\exists y \in A)x \subseteq y \lif x \in A)$. 
\end{defn}
\begin{lem}\ollabel{Valphabasicprops}
For each ordinal $\alpha$:
\begin{enumerate}
	\item\ollabel{Valphatrans} Each $V_\alpha$ is transitive.
	\item\ollabel{Valphapotent} Each $V_\alpha$ is potent.
	\item\ollabel{Valphacum} If $\gamma \in \alpha$, then $V_\gamma
	\in V_\alpha$ (and hence also $V_\gamma \subseteq V_\alpha$ by
	\olref{Valphatrans})
\end{enumerate}
\end{lem}

\begin{proof}
We prove this by a (simultaneous) transfinite induction.  For
induction, suppose that \olref{Valphatrans}--\olref{Valphacum} holds
for each ordinal $\beta < \alpha$. 

The case of $\alpha = \emptyset$ is trivial. 

Suppose $\alpha = \ordsucc{\beta}$. To show \olref{Valphacum}, if
$\gamma \in \alpha$ then $V_\gamma \subseteq V_\beta$ by hypothesis,
so $V_\gamma \in \Pow{V_\beta} = V_\alpha$. To show
\olref{Valphapotent}, suppose $A \subseteq B \in V_\alpha$ i.e., $A
\subseteq B \subseteq V_\beta$; then $A \subseteq V_\beta$ so $A \in
V_\alpha$. To show \olref{Valphatrans}, note that if $x \in A \in
V_\alpha$ we have $A \subseteq V_\beta$, so $x \in V_\beta$, so $x
\subseteq V_\beta$ as $V_\beta$ is transitive by hypothesis, and so $x
\in V_\alpha$. 

Suppose $\alpha$ is  a limit ordinal. To show \olref{Valphacum}, if
$\gamma \in \alpha$ then $\gamma \in \ordsucc{\gamma} \in \alpha$, so
that $V_\gamma \in V_{\ordsucc{\gamma}}$ by assumption, hence
$V_\gamma \in \bigcup_{\beta \in \alpha} V_\beta = V_\alpha$. To show
\olref{Valphatrans} and \olref{Valphapotent}, just observe that a
union of transitive (respectively, potent) sets is transitive
(respectively, potent). 
\end{proof}

\begin{lem}\ollabel{Valphanotref}
For each ordinal $\alpha$, $V_\alpha \notin V_\alpha$.
\end{lem}

\begin{proof}
By transfinite induction. Evidently $V_\emptyset \notin V_\emptyset$. 

If $V_{\ordsucc{\alpha}} \in V_{\ordsucc{\alpha}} = \Pow{V_\alpha}$,
then $V_{\ordsucc{\alpha}} \subseteq V_\alpha$; and since $V_\alpha
\in V_{\ordsucc{\alpha}}$ by \olref{Valphabasicprops}, we have
$V_\alpha \in V_\alpha$. Conversely: if $V_\alpha \notin V_\alpha$
then $V_{\ordsucc{\alpha}} \notin V_{\ordsucc{\alpha}}$

If $\alpha$ is a limit and $V_\alpha \in V_\alpha = \bigcup_{\beta \in
\alpha}V_\beta$, then $V_\alpha \in V_\beta$ for some $\beta \in
\alpha$; but then also $V_\beta \in V_\alpha$ so that $V_\beta \in
V_\beta$ by \olref{Valphabasicprops} (twice). Conversely, if $V_\beta
\notin V_\beta$ for all $\beta \in \alpha$, then $V_\alpha \notin
V_\alpha$.
\end{proof}

\begin{cor}
For any ordinals $\alpha, \beta$: $\alpha \in \beta$ iff $V_\alpha \in V_\beta$
\end{cor}

\begin{proof}
\Olref{Valphabasicprops} gives one direction. Conversely, suppose $V_\alpha \in V_\beta$. Then $\alpha \neq \beta$ by \olref{Valphanotref}; and $\beta \notin \alpha$, for otherwise we would have $V_\beta \in V_\alpha$ and hence $V_\beta \in V_\beta$ by \olref{Valphabasicprops} (twice), contradicting \olref{Valphanotref}. So $\alpha \in \beta$ by Trichotomy.
\end{proof}
\noindent
All of this allows us to think of each $V_\alpha$ as the $\alpha$th
stage of the hierarchy. Here is why.

Certainly our $V_\alpha$s can be thought of as being formed in an
\emph{iterative} process, for our use of ordinals tracks the notion of
iteration. Moreover, if one stage is formed before the other, i.e.,
$V_\beta \in V_\alpha$, i.e., $\beta \in \alpha$, then our process of
formation is \emph{cumulative}, since $V_\beta \subseteq V_\alpha$.
Finally, we are indeed forming \emph{all} possible collections of sets
that were available at any earlier stage, since any successor stage
$V_{\ordsucc{\alpha}}$ is the power-set of its predecessor $V_\alpha$.

In short: with $\ZFminus$, we are \emph{almost} done, in articulating
our vision of the cumulative-iterative hierarchy of sets. (Though, of
course, we still need to justify Replacement.)

% ==== OLP file: content/set-theory/spine/zf.tex ====


\olfileid{sth}{spine}{zf}
\olsection{$\Z$ and $\ZF$: A Milestone}

With Foundation, we reach another important milestone. We have
considered theories $\Zminus$ and $\ZFminus$, which we said were
certain theories ``minus'' a certain something. That certain something
is Foundation. So:

\begin{defn}
The theory $\Z$ adds Foundation to $\Zminus$. So its axioms are
Extensionality, Union, Pairs, Powersets, Infinity, Foundation, and all
instances of the Separation scheme.

The theory $\ZF$ adds Foundation to $\ZFminus$. Otherwise put, $\ZF$
adds all instances of Replacement to~$\Z$.
\end{defn}

Still, one question might have occurred to you. If Regularity is
equivalent over $\ZFminus$ to Foundation, and Regularity's
justification is clear, why bother to go around the houses, and take
Foundation as our basic axiom, rather than Regularity? 

Setting aside historical reasons (to do with who formulated what and
when), the basic reason is that Foundation can be presented without
employing the definition of the~$V_\alpha$s. That definition relied
upon all of the work of \olref[recursion]{sec}: we
needed to prove Transfinite Recursion, to show that it was justified.
But our proof of Transfinite Recursion employed \emph{Replacement}.
So, whilst Foundation and Regularity are equivalent modulo $\ZFminus$,
they are not equivalent modulo $\Zminus$. 

Indeed, the matter is more drastic than this simple remark suggests.
Though it goes well beyond this book's remit, it turns out that both
$\Zminus$ and $\Z$ are too weak to define the $V_\alpha$s. So, if you
are working only in $\Z$, then Regularity (as we have formulated it)
does not even make \emph{sense}. This is why our official axiom is
Foundation, rather than Regularity. 

From now on, we will work in $\ZF$ (unless otherwise stated), without
any further comment. 

\end{document}
